% ~7 min. The open questions hint at ongoing work.
\partframe{Wrap-up}{\phantom{Wrap-up}}

\begin{frame}
  \frametitle{Summary}
  \begin{itemize}
    \item Hylomorphisms in Rocq: \underline{axiom-free}, extractable.
    \item Laws are theorems; calculations are \code{rewrite} proofs.
    \item Termination proved once per coalgebra.
    \item Idiomatic OCaml, down to native arrays.
  \end{itemize}
  \speakernote{Containers with boolean validity give extraction without
    unsafe casts; setoids avoid functional extensionality. Dynamorphisms:
    memo tables refined to native arrays by calculation. Monoid algebras:
    one fusion law for every finite container and coalgebra.}
\end{frame}

\begin{frame}
  \frametitle{What Changed Since the Paper}
  \begin{itemize}
    \item Rocq 9; properness and termination by instance search.
    \item \code{Derive} and \code{arewrite} for calculations.
    \item Finite containers and monoid algebras.
    \item New examples: MSS, fast power, knapsack with arrays.
  \end{itemize}
  \speakernote{Also: GFix equality with Paco; mode annotations that stop
    instance search from diverging; new laws hylo\_map\_shift,
    hylo\_map\_fusion, everywhere. Skip this slide if short on time.}
\end{frame}

\begin{frame}
  \frametitle{Open Question 1: Which Intermediate Structure?}
  \[
    \begin{tikzcd}[alt={Deforestation read backwards: the hylomorphism from
        A to B is split into rana c, from A to LFix F, followed by cata a,
        from LFix F to B.}, column sep = 30mm]
      A \arrow[rr, bend left = 20, "\hylo{a}{c}"] \arrow[r, "\ana{c}"']
        \& \op{LFix}\,F \arrow[r, "\cata{a}"'] \& B
    \end{tikzcd}
  \]
  \begin{itemize}
    \item<2-> A law that splits through a \textbf{buffer}, not a tree?
    \item<3-> Fused or split: a question of \textbf{cost}.
  \end{itemize}
  \speakernote{Deforestation read right to left is fission: it puts an
    intermediate structure back, but only LFix F, a tree. A machine wants
    the intermediate in flat memory, or nowhere. Open: such a law, proved
    once, still an =e equation. Choosing fused or split depends on cost,
    which the laws do not see; could an untrusted search choose and hand
    Rocq the chain of laws as a certificate?}
\end{frame}

\begin{frame}
  \frametitle{Open Question 2: Algebras that Write}
  \begin{center}
    \begin{tikzpicture}[alt={A call tree whose leaves write to locations 0
        to 3 of a buffer. The two subtrees write disjoint sets of
        locations, 0 and 1 on the left, 2 and 3 on the right.},
      level distance = 9mm,
      level 1/.style = {sibling distance = 40mm},
      level 2/.style = {sibling distance = 20mm},
      every node/.style = {box, font = \small}]
      \node {$\cdot$}
        child {node {$\cdot$}
          child {node {$b[0] := x_0$}} child {node {$b[1] := x_1$}}}
        child {node {$\cdot$}
          child {node {$b[2] := x_2$}} child {node {$b[3] := x_3$}}};
    \end{tikzpicture}
  \end{center}
  \begin{itemize}
    \item Writes form a monoid: \code{malg\_fusion} applies.
    \item<2-> Which equivalence on buffers?
    \item<3-> Exactly once, proved per coalgebra?
  \end{itemize}
  \speakernote{Some programs return no value, they write into a buffer.
    Writes b -> b[i := x] compose into a monoid, so a writing program is a
    monoid algebra. Comparing buffers by their reads avoids functional
    extensionality, but the laws must hold up to it. Which property of the
    coalgebra makes every location written exactly once, proved once per
    coalgebra, and does it justify running subtrees at the same time? Do
    not say more than the slide.}
\end{frame}

\begin{frame}
  \frametitle{Open Question 3: Programs Without Closures}
  \begin{itemize}
    \item \code{cont} is a function: extraction gives closures.
    \item<2-> A first-order syntax, \textbf{closed under the laws}?
    \item<3-> Where do overflow and bounds conditions live?
  \end{itemize}
  \speakernote{Closures are fine for OCaml; printing for a target without
    closures needs a first-order syntax for cont (e.g. arithmetic on
    positions). fmap composes functions and fusion substitutes one into
    another: every calculated program must stay in the syntax, so that only
    the printer is trusted. The laws hold over naturals, the code runs on
    machine integers: where do the side conditions go so they do not get in
    the way of rewriting?}
\end{frame}

\begin{frame}
  \frametitle{Open Question 4: Changing the Coalgebra}
  \begin{center}
    \begin{tikzpicture}[alt={Two call trees for the same monoid algebra on
        four elements: a list-shaped tree of depth four, and a balanced
        tree of depth two. Both visit x0, x1, x2, x3 in the same order.},
      every node/.style = {box, font = \small}]
      \begin{scope}[level distance = 8mm, sibling distance = 12mm]
        \node {$\oplus$}
          child {node {$x_0$}}
          child {node {$\oplus$}
            child {node {$x_1$}}
            child {node {$\oplus$}
              child {node {$x_2$}}
              child {node {$x_3$}}}};
      \end{scope}
      \node[draw = none, fill = none, font = \Large] at (35mm, -1.2)
        {$=$};
      \begin{scope}[xshift = 70mm, level distance = 10mm,
          level 1/.style = {sibling distance = 24mm},
          level 2/.style = {sibling distance = 12mm}]
        \node {$\oplus$}
          child {node {$\oplus$} child {node {$x_0$}} child {node {$x_1$}}}
          child {node {$\oplus$} child {node {$x_2$}} child {node {$x_3$}}};
      \end{scope}
    \end{tikzpicture}
  \end{center}
  \begin{itemize}
    \item<2-> A generic law: same leaves, same order, same result?
    \item<3-> And if the order differs?
  \end{itemize}
  \speakernote{By associativity both trees agree; today this is proved per
    coalgebra, by fusing each one through the list of leaves it visits.
    Open: one law, usable as a rewrite step (split k ways, split in two
    levels). If the order differs: commutative monoids, or writes to
    distinct locations (question 2).}
\end{frame}

\begin{frame}
  \vfill
  \begin{center}
    \begin{minipage}{.6\textwidth}
      \begin{sticky}[halign = center]
        {\Huge Thank you!\par}
        \bigskip
        {\large Questions?\par}
        \medskip
        \url{https://github.com/dcastrop/rocq-hylomorphisms}
        \medskip

        {\small Based on joint work with Marco Paviotti and Michael
        Vollmer.\par}
      \end{sticky}
    \end{minipage}
  \end{center}
  \vfill
  \speakernote{Based on joint work with Marco and Michael. Questions.}
\end{frame}
